\documentclass[10pt,a4paper]{amsart}
\usepackage[margin=19mm]{geometry}
\usepackage{amsmath,amssymb,mathtools}
\usepackage[scr=rsfs,cal=euler]{mathalfa}
\usepackage{xfrac}
\usepackage[unicode,hidelinks,bookmarks=false]{hyperref}
\newcommand{\ie}{i.e.\ }
\newcommand{\cf}{cf.\ }
\newcommand{\resp}{respectively\ }
\newcommand{\Subsection}{\subsection}
\providecommand{\nobreakdash}{\nobreak\hbox{-}}
\setlength{\emergencystretch}{2em}
\multlinegap0pt
\usepackage{bm}
\usepackage{bbm}
\usepackage{dsfont}
\usepackage{xspace}
\usepackage{cancel}
\usepackage{mathtools}
\usepackage{amsthm}
\newtheorem{thm}{Theorem}
\title[Stability of compensator integrals and nonlinear operators u]{Stability of compensator integrals and nonlinear operators under canonical p-th power variation convergence (Theorem 6)}
\author{Philip Kennerberg}
\date{Source: arXiv:2605.11783v3 (CC BY 4.0)}
\begin{document}
\maketitle

\noindent\emph{Source and licence.} Stability of compensator integrals and nonlinear operators under canonical p-th power variation convergence. Version arXiv:2605.11783v3 (CC BY 4.0), distributed under the Creative Commons Attribution 4.0 International licence, as recorded in the arXiv licence metadata for this exact version and shown on the abstract page https://arxiv.org/abs/2605.11783v3. Full rights notice: licenses/arxiv-2605.11783-CC-BY-4.0.txt.\par\smallskip

\noindent\textit{Editorial scope.} Counted unit: the Thm printed as Theorem 6 of the article (source0.tex lines 4528--4625, source label \texttt{thm:weighted-cylinder-density}), delivered together with the article's own complete original proof of it (source0.tex lines 4628--5110, one \texttt{proof} environment of 483 lines). No mathematics is written, repaired or re-derived: statement and proof are the article's own text line for line. Required context retained uncounted: none. Every definition, display and label of those lines is retained unchanged. Omitted from this package and not redistributed: 1--4527, 4626--4627, 5111--13017. Nothing inside a retained excerpt is omitted or altered: every retained body is delivered exactly as the article's lines are written, so no omission receipt of any kind accompanies this package. Citations: the retained excerpts carry no citation command at all, so no bibliography entry of the article is transcribed and no reference is redistributed. Reference adaptation: none was needed and none was made; every reference inside the delivered unit resolves inside it, so no unresolved reference or citation is produced. Printed slips kept literal: no printed word, symbol, hypothesis or number of the counted statement or of its proof was altered. Layout and class: the article's own class is \texttt{article} with packages including geometry, babel, amsthm, amsmath, amssymb, natbib; the wrapper is the lane's pinned \texttt{amsart} editorial wrapper at 10pt on A4, which restates the theorem-like environments the retained text needs and adds the article's own macro definitions that the retained text needs (none), with extra wrapper packages (bm, bbm, dsfont, xspace, cancel, mathtools). The delivered numbering is the unit's own. Assets: no retained span contains a figure environment, a graphics-inclusion command, a PDF-page inclusion, a table or a TikZ picture; no image file is copied into this package, the package declares no asset dependency and no image file is redistributed with it. Licence: version arXiv:2605.11783v3 of this article is distributed under the Creative Commons Attribution 4.0 International licence, as recorded in the arXiv licence metadata for this exact version and shown on the abstract page https://arxiv.org/abs/2605.11783v3. A full rights notice is delivered at licenses/arxiv-2605.11783-CC-BY-4.0.txt. Characters: every retained excerpt is pure ASCII, so the delivered engine, XeLaTeX with its default Latin Modern text font, reports no missing character.
\begin{thm}[Characterisation of the weighted cylinder closure]
\label{thm:weighted-cylinder-density}

Let
\[
        E=\{x\in\mathbb R:b\le |x|\le a\},
        \qquad 0<b<a<\infty,
\]
and let
\[
        K(E)
        :=
        \bigl\{
        \lambda\in\mathcal M_+(E):
        \lambda(E)\le1
        \bigr\},
\]
equipped with the topology of weak convergence. Let \(H\) be a
separable Hilbert space and let \(p>1\).

Denote by \(\mathcal C_{\mathrm{cyl}}(E;H)\) the class of functions
\(P:K(E)\to H\) of the form
\[
P(\lambda)
=
\Phi
\left(
        \int_E h_1(x)\,\lambda(dx),
        \ldots,
        \int_E h_d(x)\,\lambda(dx)
\right),
\]
where
\[
        d\in\mathbb N,
        \qquad
        h_1,\ldots,h_d\in C(E),
\]
and
\[
        \Phi:\mathbb R^d\longrightarrow H
\]
is globally Lipschitz and satisfies \(\Phi(0)=0\).

For a function \(G:K(E)\to H\) satisfying \(G(0)=0\), set
\[
\lVert G\rVert_{p,*}
:=
\sup_{\substack{\lambda\in K(E)\\\lambda\ne0}}
\frac{\lVert G(\lambda)\rVert_H}
     {\lambda(E)^{1/p}}.
\]
Let
\[
\mathcal X_{p,*}(E;H)
:=
\left\{
G:K(E)\to H:
G(0)=0,\
\lVert G\rVert_{p,*}<\infty
\right\}.
\]
Equipped with \(\lVert\cdot\rVert_{p,*}\),
\(\mathcal X_{p,*}(E;H)\) is a normed vector space, and
\[
        \mathcal C_{\mathrm{cyl}}(E;H)
        \subset
        \mathcal X_{p,*}(E;H).
\]
The closure below is taken in this normed space. Then
\[
\begin{aligned}
&
\overline{\mathcal C_{\mathrm{cyl}}(E;H)}
^{\,\lVert\cdot\rVert_{p,*}}
=
\left\{
F\in C(K(E);H):
F(0)=0,\
\lim_{\delta\downarrow0}
\sup_{\substack{\lambda\in K(E)\\
                 0<\lambda(E)\le\delta}}
\frac{\lVert F(\lambda)\rVert_H}
     {\lambda(E)^{1/p}}
=0
\right\}.
\end{aligned}
\]

Moreover, for every function \(F\) in the right-hand side, the
approximating cylinder functions \(P_m\) may be chosen so that, for
every \(m\), there exists \(\delta_m>0\) such that
\[
        P_m(\lambda)=0
        \qquad\text{whenever}\qquad
        \lambda(E)\le\delta_m.
\]
\end{thm}

\begin{proof}
We first prove that every function in the set on the right-hand side
belongs to the weighted cylinder closure. Let
$
        F\in C(K(E);H)
$
satisfy \(F(0)=0\) and
\[
\lim_{\delta\downarrow0}
\sup_{\substack{\lambda\in K(E)\\
                 0<\lambda(E)\le\delta}}
\frac{\lVert F(\lambda)\rVert_H}
     {\lambda(E)^{1/p}}
=
0.
\]
Write
\[
        K:=K(E).
\]
Since \(E\) is compact, the set \(K\) is compact with respect to weak
convergence of finite measures.

For \(h\in C(E)\), define
\[
        L_h(\lambda)
        :=
        \int_E h(x)\,\lambda(dx),
        \qquad \lambda\in K.
\]
Let \(\mathfrak A\) be the unital algebra generated by the functions
\(L_h\), \(h\in C(E)\). Thus every element of \(\mathfrak A\) is a
polynomial in finitely many functions of the form \(L_h\).

The algebra \(\mathfrak A\) separates the points of \(K\). Indeed, if
\(\lambda,\widetilde\lambda\in K\) and
\[
        L_h(\lambda)=L_h(\widetilde\lambda)
        \qquad\text{for every }h\in C(E),
\]
then the uniqueness part of the Riesz representation theorem gives
$
        \lambda=\widetilde\lambda.
$
Consequently, the real Stone--Weierstrass theorem yields
\[
        \overline{\mathfrak A}^{\,\lVert\cdot\rVert_\infty}
        =
        C(K;\mathbb R).
\]

We first record the corresponding \(H\)-valued approximation
consequence. Let \(\gamma>0\). Since \(F(K)\) is compact in \(H\), there
exists a finite-dimensional subspace \(H_\gamma\subset H\) such that
\[
        \sup_{\lambda\in K}
        \operatorname{dist}
        \bigl(F(\lambda),H_\gamma\bigr)
        <
        \frac{\gamma}{2}.
\]
Let
$
        e_1,\ldots,e_N
$
be an orthonormal basis of \(H_\gamma\), and let
\[
        \Pi_\gamma:H\longrightarrow H_\gamma
\]
be the orthogonal projection. It follows that
$
        \sup_{\lambda\in K}
        \lVert F(\lambda)-\Pi_\gamma F(\lambda)\rVert_H
        <
        \frac{\gamma}{2}.
$

For \(j=1,\ldots,N\), the scalar function
\[
        \lambda
        \longmapsto
        \bigl\langle F(\lambda),e_j\bigr\rangle_H
\]
belongs to \(C(K;\mathbb R)\). By Stone--Weierstrass, there exists
\(q_j\in\mathfrak A\) such that
\[
        \sup_{\lambda\in K}
        \left|
        \bigl\langle F(\lambda),e_j\bigr\rangle_H
        -
        q_j(\lambda)
        \right|
        <
        \frac{\gamma}{2\sqrt{N}}.
\]
Define
\[
        Q(\lambda)
        :=
        \sum_{j=1}^N q_j(\lambda)e_j.
\]
Then
\[
\begin{aligned}
\lVert \Pi_\gamma F(\lambda)-Q(\lambda)\rVert_H^2
&=
\sum_{j=1}^N
\left|
\bigl\langle F(\lambda),e_j\bigr\rangle_H
-
q_j(\lambda)
\right|^2
<
\frac{\gamma^2}{4}.
\end{aligned}
\]
Therefore
$
        \sup_{\lambda\in K}
        \lVert F(\lambda)-Q(\lambda)\rVert_H
        <
        \gamma.
$ Since \(q_1,\ldots,q_N\) belong to \(\mathfrak A\), there exist
\(d\in\mathbb N\), functions
\[
        h_1,\ldots,h_d\in C(E),
\]
and a polynomial map
\[
        \Psi:\mathbb R^d\longrightarrow H_\gamma
\]
such that
\[
        Q(\lambda)
        =
        \Psi
        \left(
        \int_Eh_1(x)\,\lambda(dx),
        \ldots,
        \int_Eh_d(x)\,\lambda(dx)
        \right).
\]

We now incorporate the weight at the zero measure. Fix
\(\varepsilon>0\). By the assumed behaviour of \(F\) at zero, there
exists
$
        0<\delta<\frac12
$
such that
\[
\sup_{\substack{\lambda\in K\\
                 0<\lambda(E)\le2\delta}}
\frac{\lVert F(\lambda)\rVert_H}
     {\lambda(E)^{1/p}}
<
\frac{\varepsilon}{3}.
\]
Apply the preceding \(H\)-valued approximation result with
\[
        \gamma
        :=
        \frac{\varepsilon}{3}\delta^{1/p}.
\]
We then obtain a cylinder function \(Q\) satisfying
\[
        \sup_{\lambda\in K}
        \lVert F(\lambda)-Q(\lambda)\rVert_H
        <
        \frac{\varepsilon}{3}\delta^{1/p}.
\]

Choose a Lipschitz function
\[
        \chi_\delta:\mathbb R\longrightarrow[0,1]
\]
such that
\[
        \chi_\delta(r)=0
        \quad\text{for }r\le\delta
\]
and
\[
        \chi_\delta(r)=1
        \quad\text{for }r\ge2\delta.
\]
For example, one may take \(\chi_\delta\) to be affine on
\([\delta,2\delta]\).

Set
\[
        P(\lambda)
        :=
        \chi_\delta(\lambda(E))Q(\lambda).
\]
Since
$
        \lambda(E)=\int_E1\,\lambda(dx),
$
this is again a cylinder function. Moreover,
\[
        P(\lambda)=0
        \qquad\text{whenever}\qquad
        \lambda(E)\le\delta.
\]

We verify the weighted approximation estimate. If
\[
        0<\lambda(E)\le\delta,
\]
then \(P(\lambda)=0\), and hence
\[
\frac{\lVert F(\lambda)-P(\lambda)\rVert_H}
     {\lambda(E)^{1/p}}
=
\frac{\lVert F(\lambda)\rVert_H}
     {\lambda(E)^{1/p}}
<
\frac{\varepsilon}{3}.
\]

Suppose next that
\[
        \delta<\lambda(E)\le2\delta.
\]
Using
\[
\begin{aligned}
F(\lambda)-P(\lambda)
&=
\bigl(1-\chi_\delta(\lambda(E))\bigr)F(\lambda)
+
\chi_\delta(\lambda(E))
\bigl(F(\lambda)-Q(\lambda)\bigr),
\end{aligned}
\]
we obtain
\[
\begin{aligned}
\frac{\lVert F(\lambda)-P(\lambda)\rVert_H}
     {\lambda(E)^{1/p}}
&\le
\frac{\lVert F(\lambda)\rVert_H}
     {\lambda(E)^{1/p}}
+
\frac{\lVert F(\lambda)-Q(\lambda)\rVert_H}
     {\lambda(E)^{1/p}}
\\
&<
\frac{\varepsilon}{3}
+
\frac{
        \frac{\varepsilon}{3}\delta^{1/p}
     }{
        \lambda(E)^{1/p}
     }
\le
\frac{2\varepsilon}{3}.
\end{aligned}
\]

Finally, if
\[
        \lambda(E)>2\delta,
\]
then \(P(\lambda)=Q(\lambda)\), and therefore
\[
\begin{aligned}
\frac{\lVert F(\lambda)-P(\lambda)\rVert_H}
     {\lambda(E)^{1/p}}
&<
\frac{
        \frac{\varepsilon}{3}\delta^{1/p}
     }{
        (2\delta)^{1/p}
     }
<
\frac{\varepsilon}{3}.
\end{aligned}
\]
Combining the three cases gives
\[
\sup_{\substack{\lambda\in K\\ \lambda\ne0}}
\frac{\lVert F(\lambda)-P(\lambda)\rVert_H}
     {\lambda(E)^{1/p}}
<
\varepsilon.
\]

It remains only to verify that \(P\) can be represented using a
globally Lipschitz map. Include the constant function \(1\) among the
cylinder coordinates and define
\[
\Gamma(\lambda)
:=
\left(
        \lambda(E),
        \int_Eh_1(x)\,\lambda(dx),
        \ldots,
        \int_Eh_d(x)\,\lambda(dx)
\right).
\]
The set
\[
        \Gamma(K)\subset\mathbb R^{d+1}
\]
is compact. Choose a compactly supported Lipschitz function
\[
        \zeta:\mathbb R^{d+1}\longrightarrow[0,1]
\]
which is identically one on a neighbourhood of \(\Gamma(K)\). Define
\[
\begin{aligned}
\Phi(z_0,z_1,\ldots,z_d)
:=
\chi_\delta(z_0)\,
\zeta(z_0,z_1,\ldots,z_d)\,
\Psi(z_1,\ldots,z_d).
\end{aligned}
\]
Because \(\zeta\) has compact support and \(\Psi\) is polynomial, the
map
\[
        (z_0,z_1,\ldots,z_d)
        \longmapsto
        \zeta(z_0,z_1,\ldots,z_d)\Psi(z_1,\ldots,z_d)
\]
is bounded and globally Lipschitz. Since \(\chi_\delta\) is bounded and
globally Lipschitz, it follows that \(\Phi\) is globally Lipschitz.
Furthermore,
$
        \Phi(0)=0$
because \(\chi_\delta(0)=0\). For every \(\lambda\in K\), we have \(\zeta(\Gamma(\lambda))=1\), and
therefore
\[
\begin{aligned}
P(\lambda)
&=
\Phi
\left(
        \lambda(E),
        \int_Eh_1(x)\,\lambda(dx),
        \ldots,
        \int_Eh_d(x)\,\lambda(dx)
\right).
\end{aligned}
\]
Thus \(P\) has the required cylinder representation.

Applying the construction with \(\varepsilon=1/m\), \(m\ge1\), produces
a sequence \(P_m\in\mathcal C_{\mathrm{cyl}}(E;H)\) such that
\[
\sup_{\substack{\lambda\in K\\ \lambda\ne0}}
\frac{\lVert F(\lambda)-P_m(\lambda)\rVert_H}
     {\lambda(E)^{1/p}}
\longrightarrow0.
\]
Each \(P_m\) vanishes whenever
$
        \lambda(E)\le\delta_m
$
for the corresponding cutoff level \(\delta_m>0\). This proves that
the set on the right-hand side of the asserted identity is contained
in the weighted cylinder closure.

Conversely, suppose that
\[
        F\in
        \overline{\mathcal C_{\mathrm{cyl}}(E;H)}
        ^{\,\lVert\cdot\rVert_{p,*}}.
\]
There then exist \(P_m\in\mathcal C_{\mathrm{cyl}}(E;H)\) such that
\[
        \lVert F-P_m\rVert_{p,*}\longrightarrow0.
\]
Since \(\lambda(E)\le1\) on \(K\), we have
\[
\begin{aligned}
\sup_{\lambda\in K}
\lVert F(\lambda)-P_m(\lambda)\rVert_H
&\le
\lVert F-P_m\rVert_{p,*}
\longrightarrow0.
\end{aligned}
\]
Thus \(F\) is a uniform limit of continuous functions on \(K\), and
hence
\[
        F\in C(K;H).
\]
Moreover, \(F(0)=0\).

Fix \(m\). Write
\[
P_m(\lambda)
=
\Phi_m
\left(
        \int_E h_{m,1}(x)\,\lambda(dx),
        \ldots,
        \int_E h_{m,d_m}(x)\,\lambda(dx)
\right),
\]
and let \(L_m\) be a global Lipschitz constant for \(\Phi_m\).
Since \(\Phi_m(0)=0\),
\[
\begin{aligned}
\lVert P_m(\lambda)\rVert_H
&\le
L_m
\left(
\sum_{j=1}^{d_m}
\left|
\int_Eh_{m,j}(x)\,\lambda(dx)
\right|^2
\right)^{1/2}
\\
&\le
L_m
\left(
\sum_{j=1}^{d_m}
\lVert h_{m,j}\rVert_\infty^2
\right)^{1/2}
\lambda(E).
\end{aligned}
\]
Consequently,
\[
\lim_{\delta\downarrow0}
\sup_{\substack{\lambda\in K\\
                 0<\lambda(E)\le\delta}}
\frac{\lVert P_m(\lambda)\rVert_H}
     {\lambda(E)^{1/p}}
=
0.
\]

Let \(\varepsilon>0\). Choose \(m\) so large that
\[
        \lVert F-P_m\rVert_{p,*}
        <
        \frac{\varepsilon}{2}.
\]
For this fixed \(m\), choose \(\delta>0\) so small that
\[
\sup_{\substack{\lambda\in K\\
                 0<\lambda(E)\le\delta}}
\frac{\lVert P_m(\lambda)\rVert_H}
     {\lambda(E)^{1/p}}
<
\frac{\varepsilon}{2}.
\]
Then
\[
\begin{aligned}
&
\sup_{\substack{\lambda\in K\\
                 0<\lambda(E)\le\delta}}
\frac{\lVert F(\lambda)\rVert_H}
     {\lambda(E)^{1/p}}
\le
\lVert F-P_m\rVert_{p,*}
+
\sup_{\substack{\lambda\in K\\
                 0<\lambda(E)\le\delta}}
\frac{\lVert P_m(\lambda)\rVert_H}
     {\lambda(E)^{1/p}}
<
\varepsilon.
\end{aligned}
\]
Hence
\[
\lim_{\delta\downarrow0}
\sup_{\substack{\lambda\in K\\
                 0<\lambda(E)\le\delta}}
\frac{\lVert F(\lambda)\rVert_H}
     {\lambda(E)^{1/p}}
=
0.
\]
This proves the reverse inclusion and completes the proof.
\end{proof}

\end{document}
